% ═══════════════════════════════════════════════════════════════
%  PART VII — APRIL 8, 2026: THE PASCAL INFORMATION FUNCTOR
% ═══════════════════════════════════════════════════════════════
\clearpage
\part{The Pascal Information Functor: Every Generalization Encodes W(3,3)}

\section{Pascal's Triangle as Information-Theoretic Backbone}
\label{sec:pascal-backbone}

Pascal's triangle is not merely a combinatorial curiosity.
We demonstrate that \emph{every known generalization} of Pascal's triangle----%
$q$-deformed, hyperbolic, multinomial, fractal, $q$-Clifford, and $q$-rational---%
encodes a distinct physical aspect of the W(3,3) theory when evaluated at $q=3$.
This remarkable universality suggests that Pascal's triangle is the
\textbf{information-theoretic skeleton} of physical reality.


\section{The $q$-Pascal Triangle: Quantum Group Structure}
\label{sec:qpascal-deep}

The Gaussian binomial coefficients $\binom{n}{k}_q$ form a $q$-deformed
Pascal triangle obeying the $q$-Pascal identity:
\begin{equation}
\binom{n}{k}_q = \binom{n-1}{k}_q + q^{n-k}\binom{n-1}{k-1}_q.
\end{equation}

At $q=3$, the $q$-integers $[n]_3 = (3^n-1)/2$ are:
\begin{equation}
\boxed{
[1]_3=1,\quad [2]_3=\mu=4,\quad [3]_3=\Phithree=13,\quad
[4]_3=v=40,\quad [5]_3=(k{-}1)^2=121.
}
\end{equation}
Every $q$-integer is a W(3,3) parameter.

\begin{theorem}[$q$-Factorial and Exceptional Algebras]
\label{thm:qfactorial}
The $q$-factorials at $q=3$ give exceptional Lie algebra dimensions:
\begin{equation}
[3]_3! = 1\times 4\times 13 = 52 = \dim(F_4).
\end{equation}
Moreover, the Gaussian Pascal row~$4$ is:
\begin{equation}
\binom{4}{k}_3 = [1,\; 40,\; 130,\; 40,\; 1] = [1,\; v,\; \Phithree\Theta,\; v,\; 1].
\end{equation}
This yields neutrino mixing from pure geometry:
\begin{equation}
\sin^2\theta_{12} = \frac{\binom{4}{1}_3}{\binom{4}{2}_3} = \frac{40}{130} = \frac{\mu}{\Phithree} = \frac{4}{13} = 0.3077
\end{equation}
(measured: $0.307\pm 0.013$, exact match).
\end{theorem}


\section{The 600-Cell: Hyperbolic Pascal Meets E${}_8$}
\label{sec:600cell}

The hyperbolic Pascal simplex, constructed from the $4$-dimensional
hyperbolic mosaic $\{4,3,3,5\}$, has the \textbf{600-cell} $\{3,3,5\}$
as its vertex figure. The 600-cell's $f$-vector is:
\begin{equation}
(V,E,F,C) = (120,\; 720,\; 1200,\; 600).
\end{equation}

\begin{theorem}[600-Cell $=$ W(3,3) $\times\, q$]
\label{thm:600cell}
The 600-cell $f$-vector divided by $q=3$ recovers W(3,3):
\begin{align}
\frac{120}{q} &= 40 = v \quad\text{(vertices of W(3,3))}, \\
\frac{720}{q} &= 240 = E \quad\text{(E${}_8$ roots / total design edges)}, \\
\frac{600}{v} &= 15 = g \quad\text{(SM gauge generators)}.
\end{align}
\end{theorem}

\begin{theorem}[Bosonic-Fermionic Separation]
The 120 vertices of the 600-cell decompose as:
\begin{equation}
120 = \underbrace{24}_{\text{24-cell}\,=\,f} + \underbrace{96}_{\text{snub 24-cell}\,=\,2q s^2}
\end{equation}
where $f=24$ is the bosonic (line) count and $96 = 2qs^2$ is the
fermionic degree-of-freedom count. The 600-cell literally
separates bosons from fermions.
\end{theorem}

\begin{corollary}[Icosahedron $=$ W(3,3) Vertex Figure]
The vertex figure of the 600-cell is the icosahedron $\{3,5\}$ with:
\begin{equation}
12\text{ vertices} = k,\qquad 30\text{ edges} = 2g,\qquad 20\text{ faces} = 2\Theta.
\end{equation}
Every number of the icosahedron is a W(3,3) parameter.
\end{corollary}

Furthermore, $E_8$ folds into \emph{two} golden-ratio-scaled 600-cells:
\begin{equation}
E_8 = 240\text{ roots} = 2\times 120 = 2\times(q\cdot v),
\end{equation}
with the folding matrix organized by the 9th row of Pascal's triangle
$[1,8,28,56,70,56,28,8,1]$, which IS the grade structure of the
Clifford algebra $\Cl(8)$.


\section{The $q$-Clifford Algebra: Quantization of Gauge Structure}
\label{sec:qclifford}

The quantized Clifford algebra $\Cl_q(n,\kappa)$ of Hayashi--Kwon--Aboumrad--Scrimshaw
has dimension $(8\kappa)^n$ and decomposes into $(2\kappa)^n$ irreducible
representations of dimension $2^n$ each.

\begin{theorem}[$\Cl_q(q,\lambda)$ Dimensional Promotion]
\label{thm:qcliff-promotion}
Setting $n=q=3$ and twist $\kappa=\lambda=2$:
\begin{equation}
\boxed{
\dim\Cl_q(q,\lambda) = (8\lambda)^q = 16^3 = 4096 = 2^{12} = 2^k = \dim\Cl(k).
}
\end{equation}
The $q$-deformation with parameters $(q,\lambda)=(3,2)$
\textbf{promotes} the 6-dimensional Clifford algebra to dimension $k=12$:
\begin{equation}
\Cl(q!) = \Cl(6) \;\xrightarrow{\;\text{$q$-deform}\;}\; \Cl_q(q,\lambda) \cong \Cl(k).
\end{equation}
The number of irreps is $(2\lambda)^q = \mu^q = 64 = 2^{q!}$,
and each has dimension $2^q = 8$.
\end{theorem}

This reveals a deep structure: the classical Clifford algebra $\Cl(6)$
gives the SM gauge group ($g = \binom{6}{2} = 15$ bivectors);
its $q$-deformation $\Cl_q(3,2)$ gives the \emph{full} W(3,3) theory
at valence $k=12$.


\section{Fractal Pascal: Self-Similar Dimensions}
\label{sec:fractal-pascal}

Reducing Pascal's $d$-simplex modulo a prime $p$ produces a fractal
whose Hausdorff dimension is:
\begin{equation}
D_d(p) = \log_p\binom{p+d-1}{d}.
\end{equation}

\begin{theorem}[Fractal Dimension Sequence at $q=3$]
\label{thm:fractal-sequence}
For $d$-dimensional Pascal simplices modulo $q=3$:
\begin{center}
\begin{tabular}{@{}cccl@{}}
\toprule
$d$ & $\binom{q+d-1}{d}$ & $D_d$ & W(3,3) meaning \\
\midrule
1 & 3 $= q$ & 1 & field order \\
2 & 6 $= q!$ & $\log_3 6 \approx 1.631$ & permutation group \\
3 & 10 $= \Theta$ & $\log_3 10 \approx 2.096$ & perpendicular defect \\
4 & 15 $= g$ & $\log_3 15 \approx 2.465$ & gauge generators \\
5 & 21 & $\log_3 21 \approx 2.771$ & triangular $T_6$ \\
6 & 28 & $\log_3 28 \approx 3.033$ & $\binom{8}{2}$ \\
\bottomrule
\end{tabular}
\end{center}
\end{theorem}

\begin{theorem}[Pascal's Self-Reference]
\label{thm:pascal-self-ref}
The fractal dimensions satisfy:
\begin{equation}
\boxed{
D_d(q) = \log_q\left[\binom{q+d-1}{q-1}\right] = \log_q T_{d+1}
}
\end{equation}
where $T_n = n(n+1)/2$ is the $n$-th triangular number.
Thus \textbf{the fractal dimension of each Pascal $d$-simplex mod~$q$
is encoded in Pascal's own second diagonal}, the triangular numbers.
Pascal's triangle measures its own fractal complexity.
\end{theorem}


\section{The Multinomial Pascal $d$-Simplex}
\label{sec:multinomial}

The $d$-simplex of multinomial coefficients has level-$n$ sum $d^n$.
At $d=q=3$:
\begin{equation}
\text{Level } q \text{ of the } q\text{-simplex} = q^q = 27 = q^3 = \dim(E_6\text{ fund.}).
\end{equation}

\begin{theorem}[The $d=q!{-}1=5$ Simplex Encodes the Standard Model]
The face vector of the 5-simplex (hexateron) is Pascal row $q!=6$:
\begin{equation}
[1,\; 6,\; 15,\; 20,\; 15,\; 6,\; 1].
\end{equation}
Its entries are:
\begin{itemize}[nosep]
\item $6 = q!$ vertices (permutation group);
\item $15 = g$ edges (gauge generators / bivectors);
\item $20$ faces (icosahedron faces);
\item Total $= 2^{q!} = 64 = \mu^q$ sub-faces.
\end{itemize}
\end{theorem}


\section{The Unified Pascal Information Functor}
\label{sec:pascal-functor}

We can now state the unifying principle:

\begin{theorem}[Pascal Information Functor]
\label{thm:pascal-functor}
Every generalization of Pascal's triangle, when evaluated at $q=3$,
is a \textbf{projection} of the W(3,3) information structure:
\begin{center}
\renewcommand{\arraystretch}{1.3}
\begin{tabular}{@{}lll@{}}
\toprule
\textbf{Generalization} & \textbf{W(3,3) Aspect} & \textbf{Key Identity} \\
\midrule
Standard ($q\to 1$) & Clifford algebra $\Cl(q!)$ & $g = \binom{q!}{2}$ \\
$q$-Pascal ($q=3$) & Quantum group $\SU_q(2)$ & $[2]_3 = \mu = 4\mathrm{D}$ \\
Hyperbolic $\{4,3,3,5\}$ & 600-cell $\to E_8$ & $120/q = v$ \\
$d$-Simplex & Dimensional hierarchy & $d^q = $ level sum \\
$q$-Clifford $\Cl_q(q,\lambda)$ & Gauge quantization & $\dim = 2^k$ \\
Fractal (mod $q$) & Self-similarity & $D_d = \log_q T_{d+1}$ \\
$q$-Rationals & Farey graph & $[1/2]_3 = 1/\mu$ \\
\bottomrule
\end{tabular}
\end{center}
The three master identities connecting these projections are:
\begin{align}
\dim[\Cl_q(q,\lambda)] &= 2^k, \label{eq:master1} \\
D_d(q) &= \log_q T_{d+1}, \label{eq:master2} \\
f_i(\text{600-cell}) / q^{a_i} &= \text{W(3,3) parameter}, \label{eq:master3}
\end{align}
where $a_i \in \{1, 1, 0, 0\}$ for $(V,E,F,C)$ respectively.
\end{theorem}

The conclusion is inescapable: Pascal's triangle, far from being
a mere number pattern, is the \textbf{universal information-theoretic
structure} that generates all of physics when specialized to the
unique prime $q=3$ satisfying $(q{-}2)!=1$ and $(q{-}1)!\equiv -1\pmod{q}$.
